\section{Introduction}

Tensor decomposition provides a natural low-rank representation for multiway data and has been widely used in image restoration, signal processing, compact representation, and missing-data completion tasks \cite{Kolda2009Tensor,DeLathauwer2000MLSVD,Cichocki2015SignalProcessing,Sidiropoulos2017TSP,Liu2013TensorCompletion}. As one of the classical models, Tucker decomposition characterizes couplings across modes through a core tensor and mode-wise factor matrices, achieving a favorable balance among structural interpretability, flexibility in rank control, and algorithmic feasibility \cite{tucker1966some,de2000multilinear,Kolda2009Tensor}. However, standard Tucker decomposition is fundamentally built on a multilinear entry-generation assumption. When the observed tensor is produced by nonlinear generation, saturated responses, complex higher-order interactions, or heteroscedastic perturbations, a purely multilinear model often leads to systematic mismatch, which in turn limits its performance in reconstruction, recovery, and downstream prediction tasks.

To address this issue, a variety of nonlinear tensor decomposition methods have appeared in recent years. Existing work modifies the mapping from latent linear predictors to observed entries through link functions, generalized likelihoods, kernel methods, or Gaussian processes on the one hand, and enhances model expressivity through deep priors, neural decoders, or networkized decomposition structures on the other hand \cite{hong2020generalized,xu2011infinite,saragadam2024deeptensor,varolgunes2023nmtucker,zhai2024quadratic}. These methods show that nonlinear extensions can indeed substantially improve the ability of tensor models to adapt to complex data-generating mechanisms. At the same time, however, they are often accompanied by higher inference and training complexity, a larger hyperparameter design space, or reduced structural transparency relative to classical Tucker decomposition. It therefore remains meaningful to ask whether one can introduce controlled yet effective nonlinear expressive power into Tucker while preserving, as far as possible, its low-rank backbone, parameter efficiency, and multiway interpretability. This is precisely the starting point of this paper.

We propose Nonlinear Tucker Decomposition with Polynomial Links (NTD-PL). The model uses a shared Tucker latent tensor
\[
\cS=\cX\times_1 \mathbf{A}^{(1)}\times_2\cdots\times_N \mathbf{A}^{(N)},
\]
as the low-rank backbone, and imposes on it an explicit element-wise polynomial link
\[
\widehat{\cY}=f_{\boldsymbol{\beta}}(\cS)=\sum_{p=0}^{P}\beta_p\cS^{\odot p}.
\]
In this way, standard Tucker is naturally included as a linear special case, while higher-order nonlinearity is introduced in a controlled manner through a small number of link coefficients. Unlike nonparametric mappings based on kernels or Gaussian processes, the nonlinearity in NTD-PL is explicit, low-dimensional, and deterministic. Unlike deep nonlinear tensor models that rely on neural decoders, NTD-PL does not replace the low-rank Tucker backbone; instead, it strengthens the entry-generation mechanism in a structured way. It is therefore better understood as a model that enhances Tucker rather than replaces it.

The key point of this construction is that it is not simply an element-wise post-processing step applied to a Tucker reconstruction. Since each entry of the latent tensor $\cS$ is already generated by multilinear coupling between the core tensor and the mode factors, $\cS^{\odot p}$ corresponds to a recombination of higher-order interactions on the shared low-rank backbone rather than to an external correction independent of the low-rank structure. Precisely because of this, NTD-PL can achieve expressive power beyond purely multilinear models without explicitly introducing higher-order factor tensors, while still preserving a backbone parameterization of the same order as Tucker.

Around this model, we further characterize several of its basic properties. First, NTD-PL and standard Tucker form a clear nesting relation: when the higher-order coefficients are turned off, the model exactly degenerates to Tucker; as the polynomial degree increases, the model class expands in a controlled way with degree. Second, when the observation tensor is generated by applying an element-wise polynomial function to some Tucker latent tensor, NTD-PL can achieve exact matching at the corresponding order; when the generating function is analytic, finite-order polynomial links provide systematic approximation with explicit truncation error bounds. Third, in order to handle both full-observation tensor decomposition and tensor completion with masked observations in a unified way, we construct a unified learning objective and combine coefficient ridge updates, factor normalization, and degree continuation to obtain a practical optimization framework.

Experimental results further support the proposed modeling and theoretical analysis. In controlled synthetic experiments, NTD-PL remains consistent with Tucker on strictly linear data, and under polynomial and smooth nonlinear residuals it consistently outperforms Tucker, CP, TT, and TR as nonlinear strength increases, with performance gains aligned with the dominant order of the target function. On real hyperspectral data, NTD-PL consistently outperforms Tucker on both full-observation reconstruction and random-missing completion on the CAVE dataset under the same parameter budget. These gains appear not only in scene-wise averages but also consistently across most scenes, multiple missing rates, and local spatial regions. Further mechanism analysis shows that this advantage cannot be attributed either to simply adding more parameters or to applying a single shared polynomial calibration to the Tucker output. The main contributions of this paper are summarized as follows.

\begin{enumerate}
    \item We propose NTD-PL, a nonlinear low-rank tensor model that imposes an explicit element-wise polynomial link on a shared Tucker latent tensor.
    \item We clarify its nesting relation with Tucker, its parameter-scale characteristics, and its expressivity under polynomial and analytic generation.
    \item We provide a unified masked learning formulation so that the model can handle both full tensor decomposition and tensor completion.
    \item We validate the effectiveness and applicability boundary of the model through controlled synthetic experiments and real hyperspectral experiments.
\end{enumerate}

The remainder of the paper is organized as follows. Section~\ref{sec:related} reviews related work. Section~\ref{sec:preliminaries} introduces notation and preliminaries. Section~\ref{sec:method} presents the model, its properties, and the optimization method. Section~\ref{sec:experiments} reports the experimental results. Section~\ref{sec:conclusion} concludes the paper and discusses future work.
